% ---------------------------------------------------------------------
%  TÓM TẮT
% ---------------------------------------------------------------------
\begin{abstract}

Khi người dùng mô tả triệu chứng sức khỏe bằng ngôn ngữ tự nhiên, việc phân
tầng mức độ nguy cơ đúng đắn và kịp thời quyết định chất lượng chăm sóc ban
đầu; nếu dựa thuần vào trí tuệ nhân tạo tạo sinh, kết luận có thể thiếu nhất
quán tại các ngưỡng an toàn tính mạng. Đồ án xây dựng \textbf{hệ thống trợ
lý sàng lọc sức khỏe và phân tầng nguy cơ y tế} theo kiến trúc \textbf{Hybrid
AI}: một động cơ phân loại tất định dựa trên bộ tri thức luật y khoa đóng vai
trò quyết định, còn mô hình ngôn ngữ lớn đảm nhận trích xuất sự kiện triệu
chứng từ câu nói tự nhiên và sinh lời diễn giải lâm sàng. Bộ tri thức gồm
31 triệu chứng, 22 dấu hiệu cờ đỏ theo quy ước phân loại khẩn cấp Manchester
và 44 luật quyết định, ánh xạ ra bốn mức: cấp cứu, đi khám, tư vấn dược sĩ
và tự chăm sóc. Bắt cờ đỏ cấp cứu luôn ưu tiên tuyệt đối trong suốt quá
trình ghi nhận câu trả lời; kết quả thực nghiệm trên 26 ca bệnh mẫu cho thấy
không ca cấp cứu nào bị xếp xuống nhóm tự chăm sóc, tỷ lệ xử trí an toàn lâm
sàng đạt 88,5\% và chưa từng hạ thấp mức phân loại khi dữ liệu trích xuất do
AI sinh ra nhiễu. Lớp diễn giải sử dụng cơ chế gọi
công cụ (tool-calling) để sinh năm loại nội dung chăm sóc, theo dõi và đề
xuất sử dụng thuốc không kê đơn, được kiểm duyệt chặt về số lượng mục và
ràng buộc ngưỡng an toàn trước khi hiển thị. Hệ thống được triển khai trên
nền tảng web Laravel với công cụ quản trị tương tác dược sĩ, đảm bảo mọi
luồng an toàn lâm sàng đều có cơ chế tự động và nhân viên chịu trách nhiệm
rõ ràng.

\vspace{0.8em}
\noindent\textbf{Từ khoá:} trí tuệ nhân tạo lai, phân tầng nguy cơ y tế,
hệ chuyên gia dựa trên luật, mô hình ngôn ngữ lớn, gọi công cụ, sàng lọc triệu
chứng, an toàn lâm sàng.
\end{abstract}